\chapter{Kombinatoryka i modele skończone}
\label{ch:combinatorics}

\section{Zasady zliczania}
Zacznijmy od najprostszego problemu: chcemy obliczyć szansę
pewnego wyniku, ale najpierw musimy ustalić, ile wyników w ogóle
rozważamy. Przy rzutach kostką, losowaniu kart albo układaniu haseł
istotne są trzy pytania: czy obiekty są rozróżnialne, czy kolejność
ma znaczenie i czy dopuszczamy powtórzenia. Zmiana choćby jednej
odpowiedzi prowadzi do innego modelu.

Zasada iloczynu nie wymaga probabilistycznej niezależności.
Jest twierdzeniem o liczbie możliwych ciągów wyborów, podczas gdy
niezależność jest własnością \emph{miary} nadanej zdarzeniom.
Dla skończonych modeli te dwa pojęcia często występują obok siebie,
ale nie wolno ich utożsamiać.

\begin{definicja}
Dla liczby naturalnej $n\geq 0$ definiujemy $n!:=\prod_{j=1}^n j$, przyjmując $0!:=1$.
Dla liczb całkowitych $n\geq0$ i $0\leq k\leq n$ definiujemy
$\binom nk:=n!/[k!(n-k)!]$; dla pozostałych całkowitych $k$ przyjmujemy
$\binom nk:=0$.
\end{definicja}

\begin{stwierdzenie}[Zasady sumy i iloczynu]
Niech $A_1,\dots,A_m$ będą skończonymi zbiorami.
Jeżeli są parami rozłączne, to $|\bigcup_{i=1}^m A_i|=\sum_{i=1}^m|A_i|$.
Bez żadnego założenia o rozłączności zachodzi
$|A_1\times\cdots\times A_m|=\prod_{i=1}^m|A_i|$.
\end{stwierdzenie}
\begin{proof}
Dla rozłącznych zbiorów odwzorowanie
$\coprod_{i=1}^m A_i\to\bigcup_{i=1}^m A_i$
wysyłające element z $i$-tego składnika na ten sam element
jest bijekcją, ponieważ każdy element sumy należy do
dokładnie jednego z $A_i$. Skończony zbiór jest więc
rozbity na rozłączne części o licznościach $|A_i|$.

Dla dwóch zbiorów $A,B$ rozbijmy $A\times B$ na włókna
$\{a\}\times B$ indeksowane przez $a\in A$.
Są parami rozłączne i każde ma dokładnie $|B|$ elementów.
Pierwsza zasada daje
$|A\times B|=\sum_{a\in A}|B|=|A|\,|B|$.
Dla $m>2$ traktujemy
$A_1\times\cdots\times A_m$ jako
$(A_1\times\cdots\times A_{m-1})\times A_m$
i stosujemy powyższy argument indukcyjnie.
Ważne, że jest to dowód czysto zbiorowy i w ogóle
nie korzysta z pojęcia prawdopodobieństwa.
\end{proof}

\begin{twierdzenie}[Permutacje, wariacje i kombinacje]
Niech $S$ będzie zbiorem $n$ różnych elementów i niech $0\leq k\leq n$.
Wówczas liczba bijekcji $S\to S$ wynosi $n!$, liczba uporządkowanych
$k$-elementowych ciągów bez powtórzeń wynosi
$(n)_k:=n!/(n-k)!$, a liczba $k$-elementowych podzbiorów $S$ wynosi
$\binom nk$. Liczba ciągów długości $k$ o wyrazach w $S$, z dopuszczeniem
powtórzeń, wynosi $n^k$.
\end{twierdzenie}
\begin{proof}
Uporządkowany ciąg $(s_1,\ldots,s_k)$ bez powtórzeń
budujemy kolejno. Pierwszy wyraz ma $n$ możliwych wartości.
Po wyborze pierwszych $j-1$ wyrazów pozostało dokładnie
$n-j+1$ wartości dla $s_j$. Z zasady iloczynu liczba
ciągów jest równa
$n(n-1)\cdots(n-k+1)=n!/(n-k)!$.
Przy $k=0$ otrzymujemy jeden ciąg pusty,
a przy $k=n$ wszystkie bijekcje, czyli permutacje.

Niech $\mathcal V_k$ będzie zbiorem tych ciągów, a
$\mathcal C_k$ zbiorem $k$-elementowych podzbiorów $S$.
Odwzorowanie $\pi:\mathcal V_k\to\mathcal C_k$,
które zapomina kolejność, jest surjekcją. Dla
każdego $C\in\mathcal C_k$ zbiór $\pi^{-1}(C)$
zawiera dokładnie $k!$ uporządkowań elementów $C$.
Wszystkie włókna mają więc tę samą liczność, skąd
$|\mathcal V_k|=k!|\mathcal C_k|$ i
$|\mathcal C_k|=\binom nk$.
Dla ciągów z powtórzeniami każde miejsce ma $n$
możliwych wartości bez względu na wcześniejsze wybory,
więc ich liczba to $n^k$.
\end{proof}

\begin{przyklad}[Kod dostępu]
Kod złożony z sześciu cyfr, w którym dopuszcza się początkowe zero,
może mieć $10^6$ postaci. Jeżeli cyfry nie mogą się powtarzać,
postaci jest $(10)_6$. Jeżeli dodatkowo kod ma zaczynać się cyfrą niezerową,
a pozostałe cyfry nie mogą się powtarzać, otrzymujemy
$9\cdot(9)_5$ postaci. Warunki dotyczące pierwszej pozycji rozpatruje się
zatem przed zastosowaniem zasady iloczynu.
\end{przyklad}

\section{Powtórzenia i rozdzielanie obiektów}
Kombinacje opisują wybór podzbioru różnych elementów. Nie nadają
się jednak bezpośrednio do opisu rozdzielania identycznych przedmiotów:
zamiana miejscami dwóch identycznych żetonów niczego nie zmienia.
Odwrotnie, jeśli mamy kilka kopii jednego symbolu, permutacje
sztucznie oznakowanych kopii wielokrotnie liczą ten sam wynik.
Stąd wynikają dwie osobne techniki: dzielenie przez liczbę
niewidocznych permutacji oraz kodowanie przydziałów
za pomocą gwiazdek i kresek.

\begin{twierdzenie}[Permutacje multizbioru]
Niech $n_1,\ldots,n_r$ będą nieujemnymi liczbami całkowitymi, a
$n=n_1+\cdots+n_r$. Liczba ciągów długości $n$, w których symbol $i$
występuje dokładnie $n_i$ razy dla każdego $i=1,\ldots,r$, wynosi
$n!/(n_1!\cdots n_r!)$.
\end{twierdzenie}
\begin{proof}
Wyobraźmy sobie, że wszystkie wystąpienia symbolu $i$
otrzymują dodatkowe etykiety $(i,1),\ldots,(i,n_i)$.
Powstaje $n$ różnych obiektów, których ustawień jest $n!$.
Po usunięciu etykiet otrzymujemy ciąg z zadanymi
licznościami symboli. Odwzorowanie „usuń etykiety”
obejmuje wszystkie dopuszczalne ciągi.

Ustalmy jeden taki ciąg. Etykiety symbolu $i$
można przypisać jego $n_i$ wystąpieniom na $n_i!$
sposobów. Te operacje dla różnych symboli są niezależnymi
operacjami zliczania, więc łącznie dostajemy
$\prod_{i=1}^r n_i!$ oznakowanych ciągów.
Każdy ciąg nieoznakowany ma tyle samo reprezentacji.
Dzieląc $n!$ przez rozmiar jednego włókna, otrzymujemy
$n!/(n_1!\cdots n_r!)$.
Dla $n_i=0$ czynnik $n_i!=0!=1$ nie zmienia rachunku.
\end{proof}

\begin{twierdzenie}[Metoda gwiazdek i kresek]
Dla liczb całkowitych $n\geq0$ i $r\geq1$ liczba rozwiązań
$(x_1,\ldots,x_r)\in\Z_{\geq0}^r$ równania $x_1+\cdots+x_r=n$ wynosi
$\binom{n+r-1}{r-1}$. Jeżeli dodatkowo $x_i\geq1$ dla każdego $i$,
a $n\geq r$, liczba rozwiązań wynosi $\binom{n-1}{r-1}$.
\end{twierdzenie}
\begin{proof}
Dla rozwiązania $(x_1,\ldots,x_r)$ z $x_i\geq0$
budujemy słowo, w którym umieszczamy $x_1$ gwiazdek,
następnie kreskę, potem $x_2$ gwiazdek i tak dalej,
aż do $x_r$ gwiazdek. Powstaje dokładnie $n$ gwiazdek
i $r-1$ kresek, a więc słowo długości $n+r-1$.
Dopuszczamy kreski obok siebie i na brzegach,
co odpowiada współrzędnym równym zero.

Odwrotnie, z każdego słowa zawierającego $n$ gwiazdek
i $r-1$ kresek odzyskujemy jednoznacznie liczby $x_i$
przez policzenie gwiazdek w kolejnych blokach.
Mamy więc bijekcję między rozwiązaniami a wyborami
$r-1$ pozycji kresek spośród $n+r-1$ pozycji.
Stąd pierwszy wzór. Gdy wymagamy $x_i\geq1$, podstawiamy
$y_i=x_i-1\geq0$. Równanie staje się
$y_1+\cdots+y_r=n-r$, więc stosujemy pierwszy wzór
do $n-r$ i otrzymujemy $\binom{n-1}{r-1}$.
\end{proof}

\begin{przyklad}[Rozdzielanie jednakowych przedmiotów]
Rozdzielamy 12 identycznych żetonów pomiędzy cztery rozróżnialne osoby.
Liczba przydziałów bez dodatkowych warunków wynosi $\binom{15}{3}$.
Jeżeli każda osoba ma otrzymać przynajmniej dwa żetony, podstawienie
$y_i=x_i-2$ daje $\sum_i y_i=4$, zatem $\binom73$ przydziałów.
Wyniki nie odnoszą się do 12 \emph{rozróżnialnych} żetonów:
wtedy bez ograniczeń mamy $4^{12}$ przydziałów.
\end{przyklad}

\begin{stwierdzenie}[Tożsamość Pascala i dwumian Newtona]
Dla $n\geq1$ oraz $1\leq k\leq n-1$ zachodzi
$\binom nk=\binom{n-1}{k-1}+\binom{n-1}k$. Dla dowolnych liczb
$x,y\in\R$ i $n\geq0$ zachodzi
$(x+y)^n=\sum_{k=0}^n\binom nk x^k y^{n-k}$.
\end{stwierdzenie}
\begin{proof}
Ustalmy zbiór $S=\{1,\ldots,n\}$ i wyróżnijmy element
$n$. Każdy $k$-elementowy podzbiór $S$ albo zawiera $n$,
albo go nie zawiera. W pierwszym przypadku pozostałe
$k-1$ elementów można wybrać z $n-1$ elementów na
$\binom{n-1}{k-1}$ sposobów. W drugim wybieramy wszystkie
$k$ elementów spośród $n-1$, na $\binom{n-1}{k}$
sposobów. Obie rodziny są rozłączne i wyczerpują wszystkie
przypadki; daje to tożsamość Pascala.

Aby rozwinąć $(x+y)^n$, zapiszmy iloczyn $n$ czynników
$(x+y)$. Każdy wyraz po wymnożeniu uzyskujemy przez
wybranie z każdego czynnika albo $x$, albo $y$.
Jeśli wybraliśmy $x$ dokładnie $k$ razy, powstaje
$x^ky^{n-k}$. Można to zrobić na $\binom nk$ sposobów.
Zsumowanie wkładów dla $k=0,\ldots,n$ daje wzór Newtona.
Ten argument działa dla wszystkich $x,y\in\R$, również
zerowych i ujemnych: jest to tożsamość algebraiczna.
\end{proof}

\section{Zasada włączeń i wyłączeń}
Dodawanie liczebności zbiorów działa tylko wtedy, gdy zbiory
nie zachodzą na siebie. Jeśli element spełnia kilka warunków,
pojawia się wielokrotne liczenie. Zasada włączeń i wyłączeń
naprawia to przez kolejne odejmowanie i dodawanie przecięć.
Jej charakterystyczne znaki naprzemienne nie są sztuczną sztuczką:
wynikają z tego, że element należący do $r$ zbiorów ma być
ostatecznie policzony dokładnie raz.

\begin{twierdzenie}[Zasada włączeń i wyłączeń]
Dla skończonych zbiorów $A_1,\ldots,A_m$ zachodzi
\begin{equation}\label{eq:wlaczenia}
\left|\bigcup_{i=1}^m A_i\right|
=\sum_{\varnothing\ne I\subseteq\{1,\ldots,m\}}
 (-1)^{|I|+1}\left|\bigcap_{i\in I}A_i\right|.
\end{equation}
\end{twierdzenie}
\begin{proof}
Niech $U=\bigcup_{i=1}^m A_i$. Wystarczy pokazać, że
po prawej stronie każdy element $x\in U$ jest policzony
dokładnie raz, a każdy element spoza $U$ zero razy.
Ustalmy dowolny $x$. Jeśli nie należy do żadnego $A_i$,
nie jest zawarty w żadnym przecięciu występującym
w prawej stronie. Jego całkowity wkład jest zerowy.

Jeśli $x$ należy dokładnie do $r\geq1$ spośród zbiorów,
to pojawia się w przecięciach tych $j$-elementowych
rodzin indeksów, które wybieramy z owych $r$ zbiorów.
Takich rodzin jest $\binom rj$. Ponieważ ich znak to
$(-1)^{j+1}$, łączny wkład punktu $x$ jest
\[
\sum_{j=1}^r(-1)^{j+1}\binom rj
=1-\sum_{j=0}^r(-1)^j\binom rj
=1-(1-1)^r=1.
\]
Każdy element sumy jest więc policzony dokładnie raz.
Zsumowanie po wszystkich elementach skończonego $U$
kończy dowód. Zauważmy, że naprzemienne znaki wynikają
z potrzeby usunięcia wielokrotnego liczenia.
\end{proof}

\begin{przyklad}[Permutacje bez punktów stałych]
Niech $D_n$ oznacza liczbę permutacji zbioru $\{1,\ldots,n\}$ bez punktów
stałych. Dla każdego $i$ oznaczmy przez $A_i$ zbiór permutacji
ustalających $i$. Dla $k$ ustalonych pozycji istnieje $(n-k)!$ permutacji,
więc z~\eqref{eq:wlaczenia} wynika
\[
D_n=\sum_{k=0}^{n}(-1)^k\binom nk(n-k)!
=n!\sum_{k=0}^{n}\frac{(-1)^k}{k!}.
\]
Dla losowej jednostajnej permutacji prawdopodobieństwo braku punktów
stałych wynosi $D_n/n!$ i zmierza do $e^{-1}$.
\end{przyklad}

\begin{przyklad}[Dokładnie jedna własność]
Dla trzech zbiorów $A,B,C$ liczba elementów należących do dokładnie
jednego z nich jest równa
$|A|+|B|+|C|-2(|A\cap B|+|A\cap C|+|B\cap C|)+3|A\cap B\cap C|$.
Istotnie, element należący do dokładnie $r\in\{0,1,2,3\}$ zbiorów
wnosi odpowiednio $0,1,0,0$ do wskazanej kombinacji.
\end{przyklad}

\section{Przykłady modeli kombinatorycznych}
W przykładach będziemy za każdym razem jawnie opisywać
przestrzeń wyników. To zasadniczy etap rozwiązania: pięć kart
losowanych równocześnie tworzy podzbiór, pięć kolejnych wyciągnięć
tworzy uporządkowany ciąg, a obie przestrzenie mają inne liczebności.
Jeśli każdy ciąg jest jednakowo prawdopodobny, odpowiednio opisane
zdarzenie może mieć w obu modelach tę samą szansę; nie oznacza to
jednak, że wolno zmieniać mianownik bez odpowiedniej zmiany licznika.

\begin{przyklad}[Losowanie kart]
Z talii 52 różnych kart losujemy pięć kart bez zwracania i nie
uwzględniamy kolejności. Przestrzeń wyników ma $\binom{52}5$ elementów.
Liczba układów z dokładnie dwiema kartami w ustalonym kolorze
(w znaczeniu jednego z czterech kolorów po 13 kart) wynosi
$\binom{13}2\binom{39}3$. Liczba układów zawierających dokładnie
jedną parę jednakowych rang wynosi
$13\binom42\binom{12}3 4^3$: wybieramy rangę pary,
dwie karty tej rangi, trzy inne różne rangi oraz po jednej z czterech
kart każdej z tych rang. Dzięki wymogowi różnych rang nie liczymy
dwóch par ani trójki.
\end{przyklad}

\begin{przyklad}[Problem urodzin]
Zakładamy, że daty urodzin $n$ osób są niezależne i jednostajne na
zbiorze 365 dni, bez uwzględniania lat przestępnych. Dla $n\leq365$
prawdopodobieństwo, że wszystkie daty są różne, jest równe
$\prod_{j=0}^{n-1}(1-j/365)$. Zatem prawdopodobieństwo przynajmniej
jednej zgodności wynosi
$1-\prod_{j=0}^{n-1}(1-j/365)$.
Dla $n=23$ przekracza ono $1/2$. Jest to wniosek z modelu,
nie twierdzenie o rzeczywistym rozkładzie dat urodzin.
\end{przyklad}

\begin{przyklad}[Drogi na siatce]
Rozważmy drogę z $(0,0)$ do $(a,b)$, złożoną z kroków $(1,0)$ i $(0,1)$,
gdzie $a,b\in\N\cup\{0\}$. Każda droga składa się z $a+b$ kroków,
w tym $a$ poziomych. Dróg jest więc $\binom{a+b}{a}$.
Po narzuceniu przejścia przez $(c,d)$, gdzie $0\leq c\leq a$,
$0\leq d\leq b$, liczba dróg jest iloczynem
$\binom{c+d}c\binom{a+b-c-d}{a-c}$.
\end{przyklad}

\section{Zliczanie podziałów}
Dotychczas przypisywaliśmy obiekty do ponumerowanych miejsc.
Teraz interesuje nas samo rozbicie zbioru na grupy, bez etykiet
grup. Tak powstają liczby Stirlinga drugiego rodzaju. Ich
rekurencja obrazuje często stosowaną metodę: wyróżniamy ostatni
element i pytamy, czy tworzy samodzielną grupę, czy dołącza
do istniejącej.

\begin{definicja}
Dla $n,k\geq0$ liczba Stirlinga drugiego rodzaju $S(n,k)$ oznacza
liczbę podziałów zbioru $n$-elementowego na dokładnie $k$
niepustych, nierozróżnialnych bloków. Ustalamy $S(0,0)=1$,
$S(n,0)=0$ dla $n>0$ oraz $S(0,k)=0$ dla $k>0$.
\end{definicja}
\begin{stwierdzenie}
Dla $n,k\geq1$ zachodzi
$S(n,k)=kS(n-1,k)+S(n-1,k-1)$.
Liczba surjekcji ze zbioru $n$-elementowego na zbiór $k$-elementowy
wynosi $k!S(n,k)=\sum_{j=0}^k(-1)^j\binom kj(k-j)^n$.
\end{stwierdzenie}
\begin{proof}
Rozważmy podział zbioru $\{1,\ldots,n\}$ na $k$
niepustych, nierozróżnialnych bloków. Blok zawierający
$n$ jest albo singletonem $\{n\}$, albo zawiera inny
element. W pierwszym przypadku po usunięciu tego bloku
zostaje podział $n-1$ elementów na $k-1$ bloków,
co można zrobić na $S(n-1,k-1)$ sposobów.

W drugim przypadku usunięcie $n$ pozostawia podział
$n-1$ elementów na $k$ niepustych bloków. Aby odtworzyć
wyjściowy podział, musimy dodatkowo wskazać jeden z $k$
bloków, do którego dołączamy $n$; jest więc
$kS(n-1,k)$ możliwości. Przypadki są rozłączne,
co daje rekurencję.

Surjekcja $f:\{1,\ldots,n\}\to\{1,\ldots,k\}$
wyznacza podział na $k$ niepustych włókien
$f^{-1}(j)$; odwrotnie, każdy podział na $k$ bloków
daje $k!$ surjekcji po przypisaniu etykiet włóknom.
Stąd $k!S(n,k)$. Alternatywnie, wszystkich funkcji jest
$k^n$, a te, które nie są surjekcjami, pomijają co najmniej
jeden element przeciwdziedziny. Dla ustalonego
$j$-elementowego zbioru pomijanych wartości pozostaje
$(k-j)^n$ funkcji. Z zasady włączeń i wyłączeń wynika
$\sum_{j=0}^k(-1)^j\binom kj(k-j)^n$.
\end{proof}

\begin{przyklad}[Obsadzenie wszystkich pudełek]
Mamy $n\geq1$ rozróżnialnych kul i $m\geq1$ rozróżnialnych pudełek.
Każdą kulę umieszczamy w jednym pudełku. Wszystkich przydziałów
jest $m^n$, a przydziałów bez pustego pudełka jest $m!S(n,m)$.
Ta sama liczba, policzona zasadą włączeń i wyłączeń, wynosi
\[
\sum_{j=0}^{m}(-1)^j\binom mj(m-j)^n.
\]
Gdy wszystkie $m^n$ przydziałów są jednakowo prawdopodobne,
prawdopodobieństwo obsadzenia wszystkich pudełek wynosi
$m!S(n,m)/m^n$. Przykładowo dla $n=4$ i $m=3$
mamy $S(4,3)=6$, więc prawdopodobieństwo wynosi $36/81=4/9$.
\end{przyklad}

\begin{przyklad}[Dokładnie dwie pary w pokerze]
Losujemy pięć kart z talii 52 kart, w której występuje 13 rang
po cztery karty każdej rangi. Układ z dokładnie dwiema parami
otrzymujemy następująco: wybieramy dwie rangi par na $\binom{13}2$
sposobów, wybieramy po dwie z czterech kart każdej rangi
na $\binom42^2$ sposobów, a ostatnią kartę spośród
$11\cdot4$ kart o pozostałych rangach. Wynikowe
prawdopodobieństwo to
$\binom{13}2\binom42^2\cdot44/\binom{52}5$.
Wymóg, aby ostatnia karta miała inną rangę, zapobiega
wielokrotnemu zliczaniu i wyklucza układ z trzema
kartami jednej rangi.
\end{przyklad}

